% Standalone problem document, generated from the adjacent README.md.
% Compile with XeLaTeX. Mathematical content is maintained in Markdown.
\documentclass[11pt,a4paper]{article}
\usepackage[fontsize=10.5pt]{scrextend}
\usepackage[margin=22mm,headheight=15pt,headsep=9mm,footskip=11mm]{geometry}
\usepackage{fontspec}
\setmainfont{texgyrepagella}[Extension=.otf,UprightFont=*-regular,BoldFont=*-bold,ItalicFont=*-italic,BoldItalicFont=*-bolditalic]
\setsansfont{texgyreheros}[Extension=.otf,UprightFont=*-regular,BoldFont=*-bold,ItalicFont=*-italic,BoldItalicFont=*-bolditalic]
\setmonofont{texgyrecursor}[Extension=.otf,UprightFont=*-regular,BoldFont=*-bold,ItalicFont=*-italic,BoldItalicFont=*-bolditalic,Scale=MatchLowercase]
\usepackage{amsmath,amssymb}
\usepackage{unicode-math}
\setmathfont{texgyrepagella-math.otf}
\usepackage{xcolor,microtype,enumitem,needspace,fancyhdr}
\usepackage{bookmark}
\definecolor{ink}{HTML}{17344B}
\definecolor{accent}{HTML}{197D80}
\definecolor{muted}{HTML}{596773}
\hypersetup{colorlinks=true,linkcolor=accent,urlcolor=accent,pdftitle={TR-13 - Equality
of ranks for generic odd-order Hankel
tensors},pdfauthor={Open Problems in Numerical Linear Algebra}}
\urlstyle{same}
\setlength{\parindent}{0pt}
\setlength{\parskip}{5pt plus 1pt minus 1pt}
\linespread{1.04}
\setlength{\emergencystretch}{3em}
\widowpenalty=10000
\clubpenalty=10000
\displaywidowpenalty=10000
\allowdisplaybreaks[1]
\setlist{nosep,leftmargin=1.5em}
\providecommand{\tightlist}{\setlength{\itemsep}{0pt}\setlength{\parskip}{0pt}}
\providecommand{\pandocbounded}[1]{#1}
\pagestyle{fancy}
\fancyhf{}
\fancyhead[L]{\sffamily\footnotesize\color{muted}OPEN PROBLEMS IN NUMERICAL LINEAR ALGEBRA}
\fancyhead[R]{\sffamily\bfseries\footnotesize\color{accent}TR-13}
\fancyfoot[L]{\parbox[b]{0.9\textwidth}{\sffamily\fontsize{8}{10}\selectfont\color{muted}Editorial ratings; a bounded search cannot certify that a problem remains unsolved.\\Literature check: 2026-09-11}}
\fancyfoot[R]{\sffamily\footnotesize\color{muted}\thepage}
\renewcommand{\headrulewidth}{0.3pt}
\renewcommand{\headrule}{\hbox to\headwidth{\color{accent}\leaders\hrule height \headrulewidth\hfill}}
\makeatletter
\renewcommand{\section}{\Needspace{5\baselineskip}\@startsection{section}{1}{0pt}{12pt plus 2pt minus 2pt}{4pt}{\sffamily\bfseries\large\color{ink}}}
\renewcommand{\subsection}{\Needspace{4\baselineskip}\@startsection{subsection}{2}{0pt}{9pt plus 1pt minus 1pt}{3pt}{\sffamily\bfseries\normalsize\color{ink}}}
\makeatother
\setcounter{secnumdepth}{0}
\begin{document}
{\sffamily\small\color{muted}Tensor computations\par}
\vspace{3pt}
{\raggedright\hyphenpenalty=10000\exhyphenpenalty=10000\sffamily\bfseries\fontsize{21}{24}\selectfont\color{ink}Equality
of ranks for generic odd-order Hankel tensors\par}
\vspace{7pt}
{\sffamily\small\textbf{Difficulty:} challenging\quad\textbf{Importance:} interesting
to specialist\par}
{\sffamily\small\bfseries\color{accent}Status: Solved\par}
\vspace{4pt}
{\color{accent}\hrule height 0.8pt}
\vspace{6pt}
\textbf{Rating rationale:} Challenging because the remaining odd orders
require rank lower bounds beyond the proved even-order and cubic
mechanisms; specialist importance reflects the restriction to generic
Hankel tensors.

\subsection{Resolution --- 2026-09-11}\label{resolution-2026-09-11}

\textbf{Affirmative resolution by Matthew J. Colbrook}, Department of
Applied Mathematics and Theoretical Physics, University of Cambridge.
For every odd \(m\ge5\) and \(n\ge2\), a nonempty Zariski-open set of
complex Hankel tensors has ordinary rank, symmetric rank, ordinary
border rank, symmetric border rank and Vandermonde rank all equal to
\(\lceil(m(n-1)+1)/2\rceil\). A compressed three-slice Koszul flattening
gives the ordinary-border-rank lower bound, including arbitrary
unstructured limiting sequences; a dominant moment map supplies the
matching actual Vandermonde-rank upper bound. Exceptional Hankel tensors
and the separate all-tensors question TR-14 are not settled.

\href{https://github.com/ajt60gaibb/OpenProblemsInNLA/blob/main/references/colbrook-recovered-tensors-2026-09-11/manuscripts/TR-13.pdf}{Complete
manuscript, Theorem 1 and Sections 2--5} ·
\href{https://github.com/ajt60gaibb/OpenProblemsInNLA/blob/main/references/colbrook-recovered-tensors-2026-09-11/manuscripts/TR-13.tex}{Standalone
TeX} ·
\href{https://github.com/ajt60gaibb/OpenProblemsInNLA/blob/main/references/colbrook-recovered-tensors-2026-09-11/verification/reviews/TR-13-review.md}{Independent
complete-source PASS review} ·
\href{https://github.com/ajt60gaibb/OpenProblemsInNLA/blob/main/references/colbrook-recovered-tensors-2026-09-11/README.md}{Authorship
and provenance}.

The recovered AI-assisted proof passed independent agent review; this is
not external human peer review or formal certification. The original
target and dated source audit are retained. Difficulty, importance and
rating rationale are historical. No novelty or priority claim is made.

\subsection{Statement}\label{statement}

Fix an odd integer \(m\ge5\) and \(n\ge2\). A complex Hankel tensor
\(H\) of order \(m\) and dimension \(n\) has entries \[
H_{i_1\ldots i_m}=h_{i_1+\cdots+i_m-m},\qquad 1\le i_j\le n,
\] for \(h\in\mathbb C^{m(n-1)+1}\). Is there a nonempty Zariski-open
subset of this Hankel tensor space on which \[
R(H)=R_{\rm sym}(H)=\underline R(H)=\underline R_{\rm sym}(H)=R_V(H)?
\] Here \(R\) is the minimum number of arbitrary complex rank-one tensor
summands; \(R_{\rm sym}\) restricts summands to complex multiples of
\(v^{\otimes m}\). Each underlined border rank is the least integer
\(r\) admitting a sequence of complex tensors of the corresponding rank
at most \(r\) converging entrywise to \(H\). \(R_V\) further restricts
\(v\) to \[
v(a,b)=(a^{n-1},a^{n-2}b,\ldots,b^{n-1}),\quad (a,b)\ne(0,0).
\] Thus limits defining symmetric border rank stay in the symmetric
tensor space; limits defining ordinary border rank need not. Generically
\(R_V=\lceil(m(n-1)+1)/2\rceil\).

\subsection{Relevance}\label{relevance}

Hankel tensor decompositions encode sums of exponentials used in signal
reconstruction. The conjecture compares structured decomposition lengths
with the best unconstrained low-rank descriptions.

\subsection{References}\label{references}

\begin{enumerate}
\def\labelenumi{\arabic{enumi}.}
\tightlist
\item
  J. Nie and K. Ye, \emph{Hankel Tensor Decompositions and Ranks}, SIAM
  J. Matrix Anal. Appl. 40 (2019).
  \href{https://epubs.siam.org/doi/10.1137/18M1168285}{DOI};
  \href{https://arxiv.org/pdf/1706.03631}{author preprint}, §6, Question
  6.1 and Conjecture 6.2, p.16; Corollary 3.3 gives the generic
  Vandermonde rank.
\item
  L. Qi, \emph{Hankel Tensors: Associated Hankel Matrices and
  Vandermonde Decomposition}, 2014.
  \href{https://arxiv.org/pdf/1310.5470}{Primary preprint}, §1, equation
  (1), and §4, Theorem 3, for the structured decomposition background.
\end{enumerate}

\subsection{Status check --- 2026-09-10}\label{status-check-2026-09-10}

Rechecked \href{https://arxiv.org/pdf/1706.03631}{Nie--Ye, §6, Question
6.1 and Conjecture 6.2}, and searched the conjecture number and later
odd-order Hankel-rank literature. The general odd-order target remains
conjectural in that source; its even-order and order-three theorems are
already excluded here. No later proof or counterexample was located.
Status rests on the 2019 primary text plus this bounded search.
\end{document}
